\documentclass[UTF8,11pt]{ctexart}
\usepackage[a4paper,margin=27mm]{geometry}
\usepackage{amsmath,amssymb,amsthm,booktabs,xcolor}
\usepackage[colorlinks=true,linkcolor=blue!55!black,urlcolor=blue!55!black,citecolor=blue!55!black]{hyperref}
\hypersetup{pdftitle={用线搜索试探值重选步长：首次达到梯度阈值的两种结果},pdfauthor={}}
\newcommand{\R}{\mathbb{R}}
\newcommand{\norm}[1]{\lVert #1\rVert}
\newcommand{\ip}[2]{\langle #1,#2\rangle}
\newtheorem{theorem}{定理}
\newtheorem{proposition}{命题}
\renewcommand{\proofname}{证明}
\title{用线搜索试探值重选步长\\[3pt]\large 首次达到梯度阈值的两种结果}
\author{}
\date{2026 年 11 月 26 日}
\begin{document}
\maketitle
\begin{abstract}
非单调线搜索的试探点通过下降检验时，算法已经取得其函数值，尚未查询其梯度。
本文利用该函数值和当前梯度拟合射线上的二次模型，再决定是否缩短步长。
我们固定 ABBmin1 基线、线搜索参数和相对梯度阈值。在给定的 Hessian 上下界和精确查询模型下，
两种方法均在有限步内达标。在一族一维问题上，基线第二步恰好到达极小点，重选步长跳过了该点。
在另一个一维问题上，重选步长第二步达标，基线尚未达标。
两例均核对了初始步、BB 割线、线搜索和停止条件。若函数值与梯度分别按固定单价计费，
完整运行的查询费用也可以向两个方向翻转。最后给出仅凭射线函数值能够确定的梯度范数范围。
\end{abstract}
{\small 本稿由 AI 协助整理和校核，待作者逐项人工复核。公开稿隐去署名与联系信息。}
\section{问题和停止条件}
沿负梯度方向搜索时，试探点的函数值足以检验下降条件。
若这个点通过检验，其梯度尚未查询，刚取得的函数值还可用来拟合二次模型并重选步长。
本文考察这一修改如何影响算法首次达到给定的梯度精度。

纯 BB 方法的最坏渐近因子描述长轨迹的衰减速度 \cite{yysharp}。
首次达标的更新数还取决于轨迹在停止阈值附近经过哪些点，以及算法实际查询了哪些点的梯度。
我们固定基线、线搜索参数和阈值，构造达标顺序相反的两个问题，并分别核算函数值与梯度的费用。

设 $f\colon\R^n\to\R$ 是 $C^2$ 函数，满足
\[
  \tfrac12 I\preceq \nabla^2 f(x)\preceq4I,\qquad
  \norm{\nabla^2 f(x)-\nabla^2 f(z)}_{\mathrm{op}}\le\norm{x-z}.
\]
查询函数值和梯度分别计数，查询结果在数学模型中精确。写 $g_k=\nabla f(x_k)$，
$G_k=\norm{g_k}$。初次查询 $f(x_0),g_0$ 后固定 $\theta=2^{-20}G_0$。
当且仅当某个已接受的 $x_k$ 满足 $G_k\le\theta$ 时停止。
$G_0=0$ 时在初始点停止。计数 $N$ 从初始点起只计算接受的更新。

\section{两种方法使用的步长与查询}
令 $F_k=\max_{\max(0,k-9)\le i\le k}f(x_i)$。这与 Grippo、Lampariello 和 Lucidi
使用近期函数值最大值的准则同型 \cite{gll}。两种方法均从 $a=1/4$ 开始第一步。
此后按 ABBmin1 的长短步选择思路 \cite{abbmin}，用本方法实际接受的相邻点计算 $s_k=x_k-x_{k-1}$ 和 $y_k=g_k-g_{k-1}$，
再计算两个 BB 步长 \cite{bb}
\[
  b_{1,k}=\frac{s_k^\top s_k}{s_k^\top y_k},\qquad
  b_{2,k}=\frac{s_k^\top y_k}{y_k^\top y_k}.
\]
若 $b_{2,k}/b_{1,k}<4/5$，取最近至多十个 $b_{2,i}$（含当前项）的最小值，否则取 $b_{1,k}$。
把提议步长限制在 $[1/4,2]$。每次只更新一次 BB 历史。逐次将试探步长减半，直到 $p=x_k-ag_k$ 首次满足
\begin{equation}\label{eq:gll}
     f(p)\le F_k-\tfrac14 aG_k^2.
\end{equation}
减半过程中只查询试探点的函数值。记基线为 $R$。它在 $p$ 查询梯度，接受 $p$，
然后检查梯度停止条件。

记使用试探值重选步长的方法为 $C$。第一步与 $R$ 相同。此后，在得到满足
\eqref{eq:gll} 的 $f(p)$ 后、查询 $p$ 的梯度前，计算
\begin{equation}\label{eq:model}
  \beta=\frac{2(f(p)-f(x_k)+aG_k^2)}{a^2G_k^2}.
\end{equation}
这使射线上的模型 $m(t)=f(x_k)-tG_k^2+\frac12\beta t^2G_k^2$ 在 $t=a$ 与实测函数值吻合。
若 $a\beta\le1$，仍选 $p$。若 $a\beta>1$，取 $b=1/\beta$ 和 $q=x_k-bg_k$。
先取得 $f(q)$ 并检验 \eqref{eq:gll} 中相应的 $b$，再取得 $q$ 的梯度。
只有选中的点进入已接受点的函数值窗口与割线历史。特别地，跳过 $p$ 时保留已付出的
$f(p)$ 查询，但并不知道 $p$ 的梯度。若获得的 $f(q)$ 无法通过检验，实现中应在
梯度查询和接受 $q$ 前停止。本节以下的精确模型给出这项检验成立的充分条件。

\section{线搜索条件与有限步达标}
\begin{theorem}\label{thm:finite}
对于上述函数类，两种方法的实际步长 $\alpha$ 均满足 $3/16<\alpha\le2$，
每一步的初次试探和回退合计最多查询四个线搜索试探点。
在精确查询模型中，重选的 $q$ 满足原线搜索条件。
令 $f_*=\min f$，$D_k=\max_{\max(0,k-9)\le i\le k}(f(x_i)-f_*)$。
对已生成到第 $k+10$ 个接受点的轨迹，有 $D_{k+10}\le (448/451)D_k$。
因此两种方法各自在有限步内达到由初始梯度确定的固定阈值。
\end{theorem}
\begin{proof}
令 $x=x_k$、$G=G_k>0$、$u=g_k/G$。沿 $x-ta g_k$ 作两次积分，得到
\[
 \beta=2\int_0^1(1-t)u^\top \nabla^2 f(x-ta g_k)u\,dt\in[1/2,4].
\]
重选发生时 $0<b<a$。设 $\lambda=b/a$，则 $q=(1-\lambda)x+\lambda p$，
强凸性给出
\begin{align*}
 f(q)&\le (1-\lambda)f(x)+\lambda f(p)
             -\tfrac14\lambda(1-\lambda)a^2G^2\\
 &\le F_k-\tfrac14 bG^2
             -(1-\lambda)(F_k-f(x))-\tfrac14 b(a-b)G^2.
\end{align*}
此处余项均非正，故 $q$ 严格满足原窗口的检验。

对实际接受的非零位移 $s$，梯度差为 $y=As$，其中 $A$ 是线段上的平均 Hessian。
$A$ 的谱在 $[1/2,4]$ 中，所以 $s^\top y>0$、$y^\top y>0$。
由 $s^\top A s$ 的界，以及 $\tfrac12 A\preceq A^2\preceq 4A$，两个 BB 步长也都落在 $[1/4,2]$。
从 $L=4$ 的光滑性得
\[
 f(x-ag_k)\le f(x)-a(1-2a)G^2.
\]
$a\le3/8$ 必通过 \eqref{eq:gll}。因此从至多 $2$ 开始至多减半三次。
若确实减半，上一次试探 $2a$ 被拒绝，故 $a>3/16$；未减半时 $a\ge1/4$。
重选的 $b\in[1/4,2]$ 仍满足所需步长界。

写 $h_k=f(x_k)-f_*$。强凸性与光滑性分别给出 $h_k\le G_k^2$ 和 $G_k^2\le8h_k$。
前者对 $f$ 的二次下界取最小值得到，后者对 $x_k-g_k/4$ 应用下降估计得到。
线搜索随即给出 $h_{k+1}\le D_k-(3/64)G_k^2$。同时光滑性给出
\[
 h_{k+1}\le(1-\alpha+2\alpha^2)G_k^2\le7G_k^2.
\]
合并两式即 $h_{k+1}\le(448/451)D_k$。滚动最大值 $D_k$ 非增；
十次更新后，窗口中的每个值都已在该界内，故有定理中的收缩。
若尚未停止，经过 $10m$ 次更新有
$G_{10m}^2\le 8(448/451)^m h_0\le8(448/451)^mG_0^2$。
右侧最终小于 $\theta^2$。
\end{proof}

\section{首次达标可以提前，也可以推迟}
\begin{theorem}\label{thm:opposite}
在同一组参数和停止条件下，存在合法的一维问题使 $N_R=2<N_C$，
也存在合法的一维问题使 $N_C=2<N_R$。
\end{theorem}
\subsection{重选步长跳过达标点}
给定 $r\in[1/2,1]$，置 $c=2-r/2$。规定二阶导数在 $[0,r]$ 为 $c+x$，
在 $[r,2r]$ 为 $2+3r/2-x$，区间外为 $c$。取 $g(0)=f(0)=0$ 逐次积分。
二阶导数连续，分段斜率绝对值至多 $1$，取值在 $[3/2,5/2]$。
在 $0\le x\le r$，
\[
 g(x)=cx+x^2/2,\qquad f(x)=cx^2/2+x^3/6 .
\]
由于二阶导数在 $[0,2r]$ 上关于 $r$ 对称，它的一阶矩等于总质量乘 $r$。
由此 $g(2r)=4r$、$f(2r)=4r^2$。另外 $g(r)=2r$、$f(r)=r^2-r^3/12$。

从 $x_0=2r$ 开始，两种方法以 $1/4$ 走到 $x_1=r$。
第一次更新的线搜索右端是 $3r^2$，因此 $f(r)$ 通过检验。初始点和 $r$ 均未达标。
此时 $s=-r$、$y=-2r$，两个 BB 步长均为 $1/2$。于是 $a=1/2$，$p=0$，
第一次试探即通过，且 $g(p)=0$。基线第二步终止。
方法 $C$ 尚未取得 $g(p)$，从 \eqref{eq:model} 得到
\[
  \beta=2+r/6,\qquad a\beta=1+r/12>1,\qquad
  q=\frac{r^2}{12+r}.
\]
$q$ 通过线搜索，但 $q/r\ge1/26$，故 $g(q)\ge(3/2)q\ge3r/52>4r\,2^{-20}=\theta$。
因此其第二个接受点尚未达标。定理~\ref{thm:finite}保证后续有限达标。

\subsection{重选步长先达到阈值}
在 $|x|\le1/2$ 上取 $f(x)=x^2-x^4/12$。在外侧，以 $w=|x|-1/2$ 写
$f(x)=47/192+(23/24)w+(7/8)w^2$。边界处的函数值、一阶导数与二阶导数吻合。
二阶导数介于 $7/4$ 与 $2$，且全局 $1$-Lipschitz。取 $u=x_0=2^{-8}$。
第一步走到 $r=u/2+u^3/12$，实际梯度 $h=g(r)>0$，并有
\[
    u/2<r<3u/4,\quad h/G_0>1/2,\quad
    \theta=u^3(2-u^2/3)/16 .
\]
第一步按定理~\ref{thm:finite}的下降估计通过。用 $g(x)=2x-x^3/3$ 直接因式分解得
$g(r)-g(u)=c(r-u)$，其中 $c=2-(u^2+ur+r^2)/3\in(3/2,2)$。
两个 BB 步长均为 $a=1/c$。第二次更新的首次试探为
\[
     p=-d,\qquad d=\frac{ru(u+r)}{3c},\qquad 0<d<r,\quad d<u^3/3 .
\]
$f(u)>u^2/2$，$h<3u/2$，$a<2/3$，所以线搜索右端
$f(u)-ah^2/4>u^2/8$。另一方面 $f(p)<d^2<u^6/9<u^2/8$，试探点通过检验。

令 $\ell=ah=r+d$。把四次多项式在 $r$ 展开，
\[
 f(r-\ell)-f(r)+\ell h
 = (1-r^2/2)\ell^2+(r/3)\ell^3-\ell^4/12 .
\]
由此得到
\[
 \beta=2-r^2/2+rd/3-d^2/6,\qquad
 \beta-c=\frac{2u^2+2ur-r^2+2rd-d^2}{6}>0 .
\]
因此 $b<a$。由 $q=r-h/\beta$ 得
\[
    q=-e,\qquad e=\frac{r(r-d)^2}{6\beta}.
\]
这里 $e>0$，由 $b<a$ 得 $p<q<0$，且 $e<r^3/9<3u^3/64$。线搜索按定理~\ref{thm:finite}通过，
$|g(q)|=2e-e^3/3<3u^3/32<\theta$，所以 $N_C=2$。
对基线，令 $\rho=r/u>1/2$。由
\[
 \frac{d}{\theta}=\frac{16\rho(1+\rho)}{3c(2-u^2/3)}>1
\]
以及 $|g(-d)|=2d-d^3/3>d$，可知第二步的基线尚未达标。
两个例子只需核查前两次更新。尚未在第二步停止的方法也会在有限步内达标。这一点由定理~\ref{thm:finite}给出，
无须外推两条轨迹在分开后的割线关系。

\section{射线信息和查询成本}
\begin{proposition}\label{prop:info}
考虑 $n\ge2$ 的二次模型，其 Hessian $A$ 满足 $\mu I\preceq A\preceq LI$，其中 $0<\mu\le L$。
给定当前点的值、非零梯度 $g$ 以及沿 $-g$ 射线的全部函数值，设该射线的二次曲率为
$\beta=u^\top Au$，其中 $u=g/G$。令 $q=x-g/\beta$。与这些信息相容的梯度范数 $\norm{\nabla f(q)}$ 构成区间
\[
  \left[0,\frac{G}{\beta}\sqrt{(\beta-\mu)(L-\beta)}\right].
\]
\end{proposition}
\begin{proof}
写 $Au=\beta u+w$，其中 $w\perp u$。二次形下 $\nabla f(q)=-Gw/\beta$。
对半正定矩阵 $(A-\mu I)(LI-A)$ 作用于 $u$ 的二次型展开，得到
$\norm{w}^2\le(\beta-\mu)(L-\beta)$。
反过来，取任一单位向量 $v\perp u$，在 $\operatorname{span}\{u,v\}$ 上用矩阵
\[
 \begin{pmatrix}\beta&t\\t&L+\mu-\beta\end{pmatrix},
 \qquad 0\le t\le \sqrt{(\beta-\mu)(L-\beta)}.
\]
这个矩阵减去 $\mu I$ 与从 $LI$ 减去该矩阵所得的两个矩阵均半正定：
对角元非负，行列式均为 $(\beta-\mu)(L-\beta)-t^2$。
在正交补上取 $\mu I$。对每个 $t$ 固定相同的 $f(x),g$，则沿射线的全部值相同，
并且梯度范数等于 $Gt/\beta$，取遍所述区间。
\end{proof}

上述上界是经典 Bhatia--Davis 方差界的谱形式 \cite{bd}。若上端严格大于 $\theta$，同样的射线信息可以对应两个二次模型，分别使 $q$ 达标和未达标。这个信息界只涉及给定点及其射线。两条从冷启动运行的历史还可能提供横向信息。
对于算法比较，记 $J$ 为线搜索回退时被拒绝的试探总数，$M$ 为发生重选的次数，$N$ 为首次达标的更新数。
在逐次分别查询的模型中，含初始点的查询次数为
\[
    \# f=1+N+J+M,\qquad \# g=1+N .
\]
基线 $R$ 的 $M=0$。在两个例子的前两次更新中，$C$ 请求 $4$ 次函数值和 $3$ 次梯度，
$R$ 请求 $3$ 次函数值和 $3$ 次梯度。仅凭达标与否，就能比较完整运行的分项请求数。

\begin{proposition}\label{prop:cost}
在函数值、梯度分别查询且按固定单价计费的模型中，设函数值单价 $c_f\ge0$，
梯度单价 $c_g>0$。第一例中，$C$ 的完整查询费用至少比 $R$ 多 $2c_f+c_g$。
第二例中，$C$ 的完整查询费用至少比 $R$ 少 $c_g$。
\end{proposition}
\begin{proof}
第一例中 $R$ 在第二步结束，总计 $(\#f,\#g)=(3,3)$。
$C$ 至少接受三次更新，第二步已经重选一次，故总计至少 $(5,4)$。
第二例中 $C$ 在第二步结束，总计 $(4,3)$。$R$ 至少接受三次更新，故总计至少 $(4,4)$。
以上下界逐项乘单价即得。
\end{proof}
实际运行还需要支付点构造、有限精度认证与向量运算，并可能以缓存或联合查询改变单个请求的代价。
因此这个推论给出分项查询费用的严格比较，实测时间仍需要独立验证。

\section{与已有工作的关系}
BB 的长短步与非单调搜索已有成熟理论 \cite{bb,gll,fh}；射线上用两个函数值和一个导数拟合
二次模型也是常见插值。Xie、Sun 与 Yuan 明确给出了同一插值步长 \cite[式 (16)]{xsy}。
随后关于三维终止步长的研究在附录算法 3 中把该公式用于线搜索拒绝后的重试，并要求插值步长落在
原步长的 $0.1$ 至 $0.9$ 倍之间 \cite{three}。本文在试探已经通过线搜索后才考察该公式，
并在取得试探点梯度之前决定是否改用内点。这一差别决定了基线可能已经到达梯度阈值、而重选方法
仍继续的情形。AdaBB 则由无需同类回溯的步长约束建立另一种收敛保证 \cite{adabb}。
两种方法采用相同的 GLL 窗口、初始步长、BB 选择规则和梯度阈值。
轨迹分开后，各自只用实际接受的点更新历史。两个一维构造说明重选步长既可能提前、也可能推迟首次达标。
命题~\ref{prop:cost}给出了分别计价模型下的两种严格费用比较。

实际运行时间还要考虑函数值与梯度是否联合计算、能否复用缓存，以及数值认证的成本。
有限精度实现尤其需要处理 $a\beta=1$、线搜索等号和梯度阈值附近的比较。
命题~\ref{prop:info}只固定当前点的数据和射线值。利用完整迭代历史的策略仍需单独分析。

\begin{thebibliography}{9}
\bibitem{bb} J. Barzilai and J. M. Borwein. Two-point step size gradient methods.
\emph{IMA Journal of Numerical Analysis}, 8 (1988), 141--148.
\href{https://doi.org/10.1093/imanum/8.1.141}{DOI}.
\bibitem{gll} L. Grippo, F. Lampariello and S. Lucidi. A nonmonotone line search technique for Newton's method.
\emph{SIAM Journal on Numerical Analysis}, 23 (1986), 707--716.
\href{https://doi.org/10.1137/0723046}{DOI}.
\bibitem{abbmin} G. Frassoldati, L. Zanni and G. Zanghirati. New adaptive stepsize selections in gradient methods.
\emph{Journal of Industrial and Management Optimization}, 4 (2008), 299--312.
\href{https://doi.org/10.3934/jimo.2008.4.299}{DOI}.
\bibitem{bd} R. Bhatia and C. Davis. A better bound on the variance.
\emph{The American Mathematical Monthly}, 107 (2000), 353--357.
\href{https://doi.org/10.1080/00029890.2000.12005203}{DOI}.
\bibitem{fh} Y.-H. Dai, W. W. Hager, K. Schittkowski and H. Zhang. The cyclic Barzilai--Borwein method for unconstrained optimization.
\emph{IMA Journal of Numerical Analysis}, 26 (2006), 604--627.
\href{https://doi.org/10.1093/imanum/drl006}{DOI}.
\bibitem{yysharp} S. Yang and Y.-X. Yuan. The sharp worst-case asymptotic rate of the Barzilai--Borwein method in
$\R^d$ and Hilbert spaces. \href{https://arxiv.org/abs/2608.07839}{arXiv:2608.07839}, 2026.
\bibitem{adabb} D. Zhou, S. Ma and J. Yang. AdaBB: adaptive Barzilai--Borwein method for convex optimization.
\href{https://arxiv.org/abs/2401.08024}{arXiv:2401.08024}.
\bibitem{xsy} Y. Xie, C. Sun and Y.-X. Yuan. Adaptive cyclic gradient methods with interpolation.
\emph{Computational Optimization and Applications}, 92 (2025), 301--325.
\href{https://doi.org/10.1007/s10589-025-00691-y}{DOI}.
\bibitem{three} Y. Xie, J.-P. Liu, C. Sun and Y.-X. Yuan. New gradient methods with 3 dimensional quadratic termination.
\href{https://arxiv.org/abs/2602.13141}{arXiv:2602.13141}, 2026.
\end{thebibliography}
\end{document}
